\begin{helper}[Original excursion and final-departure order]
In the notation of the preceding construction, let $\beta,v,z,P$ be any
data with $v\in Z$, with $\beta$ missing at $v$, free, and distinct from
$\gamma(r)$, and with $P$ a full simple $(\delta,\beta)$ path from $v$ to
$z$ outside $V(T)$.  There is a $\beta$-colored edge $h$ with exact ends
$u^1_\delta,u^2_\delta$.  In one of its two orientations, $P$ contains
the five consecutive incidence nodes
\[
 f^1_\delta,u^1_\delta,h,u^2_\delta,f^2_\delta.
\]
This excursion occurs strictly before the last tree vertex on $P$; that
last vertex is $v^0_\delta$, its immediately following node is
$f^0_\delta$, and every later vertex-node is outside $V(T)$.  The claim
includes exact finite indices for all five consecutive nodes and for the
last vertex and following edge.
\end{helper}
\begin{proof}
Put $X=V(T)$, put $r=\gammaIndex(\delta)$, and put
$Z=X\setminus\{u^1_{\gamma(r)},u^2_{\gamma(r)}\}$.  Since $\beta$ is
missing at $v\in Z$ and $\beta\ne\gamma(r)$,
\noderef{ScheideReplacementInternalMissing} supplies an edge identity
$h$ such that
\[
 \varphi(h)=\beta,
 \qquad \operatorname{ends}(h)=
       \{u^1_\delta,u^2_\delta\}.                         \tag{1}
\]
By \noderef{ScheideReplacementData}, for $i=1,2$ the edge
$f^i_\delta$ has color $\delta$, its tree end $v^i_\delta$ lies in
$Z$, and its other end $u^i_\delta$ lies outside $X$.  Thus the
definitions \noderef{BoundaryEdges} and \noderef{BoundaryColorEdges}
put both identities in $\partial_{\varphi,\delta}(X)$.  Applying
\noderef{ScheideDefectiveChainBoundaryEdges} to the supplied full path
shows that the two edge-nodes $f^1_\delta,f^2_\delta$ occur on $P$.
The adjacencies through (1) then put the nodes
$u^1_\delta,h,u^2_\delta$ in the same component.  Since
\noderef{ScheideDefectiveChainPath} identifies the support of $P$ with
that entire component, all five named nodes occur on $P$.

We next justify their order, including the point that the component is
not a cycle.  The globally maximum field of
\noderef{ScheideReducedConfiguration} contains a proper partial coloring.
Unpacking \noderef{IsProperEdgeColoring}, two distinct edges of color
$\delta$ incident with one vertex would be adjacent in the sense of
\noderef{AdjacentEdges} and hence are impossible.  Moreover, by the
definitions \noderef{MissingColors} and \noderef{PresentColors}, no
$\beta$-colored edge is incident with the starting vertex $v$.  The walk
is nontrivial because $v\in X$ whereas its terminal vertex $z$ is outside
$X$.  Its first step therefore gives one neighbor of $v$ in the incidence
graph, and the preceding two observations show that this is its unique
neighbor.  Here \noderef{TwoColorIncidenceGraph} is used literally: a
neighbor of the vertex-node $v$ is an individual incident edge identity
whose color is $\delta$ or $\beta$.

Every incidence node has degree at most two by
\noderef{TwoColorIncidenceDegreeAtMostTwo}; that degree bound itself is
the consequence of the endpoint-pair definition in
\noderef{LooplessMultigraph} and of
\noderef{AdjacentEdges} together with
\noderef{IsProperEdgeColoring}.  We claim that whenever two nodes of the
component are adjacent, their occurrences on $P$ are consecutive.  Let
their indices be $i<j$.  If $i=0$, uniqueness of the neighbor of $v$
forces $j=1$.  If $i>0$ and $j>i+1$, the predecessor of the node at $i$,
its successor, and the node at $j$ are three distinct neighbors: they are
neighbors because $P$ is a walk, and they are distinct because $P$ is
simple.  This contradicts the degree-two bound.  Thus $j=i+1$, proving
the claim for either orientation.

The exact endpoint equations give the four incidence adjacencies
\[
 f^1_\delta-u^1_\delta-h-u^2_\delta-f^2_\delta.           \tag{2}
\]
The colors show $f^1_\delta\ne h\ne f^2_\delta$: indeed $\delta$ is not
missing at $v$, since closedness of $X$ from
\noderef{ScheideGlobalMaximumOrder} and the $\delta$-boundary edge
$f^1_\delta$ exclude that, whereas $\beta$ is missing at $v$.
Consequently $\delta\ne\beta$.  Also
$u^1_\delta\ne u^2_\delta$, since otherwise (1) would be a loop,
contrary to \noderef{LooplessMultigraph}.  Apply the preceding adjacency
claim successively to (2).  A reversal at an intermediate step would
repeat respectively $f^1_\delta$, one of the two outside vertices, or
$h$, contradicting simplicity and these inequalities.  Hence all four
steps point in the same direction.  This produces indices
$i_0<i_1<i_2<i_3<i_4$, consecutive one by one, with the five values in
the displayed order or in its reverse.

It remains to identify the last departure.  Among the finitely many
positions of $P$, choose the largest index $q_v$ whose node is a
vertex-node $w$ in $X$.  This set of positions is nonempty because it
contains the starting position.  Also $q_v<|P|$, since the terminal
vertex $z$ is outside $X$.  By bipartiteness in
\noderef{TwoColorIncidenceGraph}, the node at $q_f=q_v+1$ is an edge-node,
say $f$, and the following node is a vertex-node $x$.  Maximality of
$q_v$ gives $x\notin X$.  The two walk adjacencies say that $f$ is
incident with $w$ and $x$; simplicity gives $w\ne x$, and the unordered
endpoint pair in \noderef{LooplessMultigraph} is therefore exactly
$\{w,x\}$.  Hence \noderef{BoundaryEdges} makes $f$ a boundary identity
of $X$.  Its color is $\delta$ or $\beta$ by
\noderef{TwoColorIncidenceGraph}.  The latter is impossible: $\beta$ is
missing at $v\in X$, while closedness from
\noderef{ScheideGlobalMaximumOrder}, interpreted by
\noderef{ColoredVertexSetClosed}, forbids a $\beta$-colored boundary
edge.  Thus $\varphi(f)=\delta$.

The data $v,\beta,\delta,z,P,q_v$ and the defining maximality of $q_v$
satisfy every clause of \noderef{ScheideExitVertex} for $w$; hence $w$
belongs to \noderef{ScheideExitVertices}.  By
\noderef{ScheideStructuralExitsInBase}, $w=C.y(a)$ for some $a<p$.
We now verify the omitted implication ``base vertex, hence in $Z$''.
The ordered layout in \noderef{ScheideReducedConfiguration}, through
its inherited \noderef{ScheideStructuralConfiguration} data, places
$C.y(a)$ at tree position $a<p$, while the two deleted block vertices
occupy positions $p+2r$ and $p+2r+1$.  These positions are distinct, and
the vertex map is injective by \noderef{TreeSequence}.  Therefore $w$
is neither deleted block vertex, so $w\in Z$.

We have now shown that $f$ is a $\delta$-colored boundary identity of
$Z$, with the incident vertex $x$ outside the larger set $X$.  The
exact-three clause of \noderef{ScheideReplacementData} gives
\[
                    f\in\{f^0_\delta,f^1_\delta,f^2_\delta\}.       \tag{3}
\]
The excursion (2) returns to $X$ after its last displayed edge.  More
explicitly, in the first orientation the node following
$f^2_\delta$ must be $v^2_\delta\in Z$: its other adjacent node is
$u^2_\delta$, its exact ends are
$\{v^2_\delta,u^2_\delta\}$, and simplicity forbids returning to the
previous node.  In the reverse orientation the same argument after
$f^1_\delta$ gives $v^1_\delta\in Z$.  Thus an $X$-vertex occurs at
index $i_4+1$, and maximality of $q_v$ yields $i_4<q_v$.  Since
$q_f=q_v+1$, simplicity of $P$ now shows that its edge-node $f$ cannot
equal either already occurring identity $f^1_\delta$ or
$f^2_\delta$.  Equation (3) therefore forces $f=f^0_\delta$.

Finally, \noderef{ScheideDefectiveSetup} gives
$\operatorname{ends}(f^0_\delta)=\{v^0_\delta,u^0_\delta\}$ with
$u^0_\delta\notin X$.  Since its endpoint $w$ lies in $X$, necessarily
$w=v^0_\delta$.  The chosen indices satisfy the four successor
equalities, $i_4<q_v$, and $q_f=q_v+1$; their path values are exactly the
five nodes of (2), followed at the last departure by
$v^0_\delta,f^0_\delta$.  The maximal definition of $q_v$ says directly
that every later vertex-node lies outside $X$, which is the remaining
assertion.
\end{proof}
